\section{Visualization, ablation, and modes of error}
\seclabel{experiments}
%The CNN works well, but what did it learn, what design choices are critical, and what are its failure modes?

\subsection{Visualizing learned features}
\seclabel{vis}

\begin{figure*}[t!]
\centering
%cat-x3-y2-c248-p1.pdf
%dog-x2-y2-c142-p1.pdf
%dog-x3-y2-c142-p1.pdf
%dots-x3-y3-c204-p1.pdf
%person-x2-y3-c132-p1.pdf
%red-x6-y3-c146-p1.pdf
%sheep-x2-y3-c117-p1.pdf
%text-x5-y2-c163-p1.pdf
%windows-x3-y3-c161-p1.pdf
\includegraphics[width=\linewidth,clip=true,trim=0 0 0.2em 0]{figures/feature-explorer/person-x2-y3-c132-p1.pdf}\\
\includegraphics[width=\linewidth,clip=true,trim=0 0 0.2em 0]{figures/feature-explorer/dots-x3-y3-c204-p1.pdf}\\
\includegraphics[width=\linewidth,clip=true,trim=0 0 0.2em 0]{figures/feature-explorer/red-x6-y3-c146-p1.pdf}\\
\includegraphics[width=\linewidth,clip=true,trim=0 0 0.2em 0]{figures/feature-explorer/text-x5-y2-c163-p1.pdf}\\
\includegraphics[width=\linewidth,clip=true,trim=0 0 0.2em 0]{figures/feature-explorer/windows-x3-y3-c161-p1.pdf}\\
\includegraphics[width=\linewidth,clip=true,trim=0 0 0.2em 0]{figures/feature-explorer/specular-x3-y3-c34-p1.pdf}
\caption{\textbf{Top regions for six \pool{5} units.} 
Receptive fields and activation values are drawn in white. 
Some units are aligned to concepts, such as people (row 1) or text (4).
Other units capture texture and material properties, such as dot arrays (2) and specular reflections (6).
}
\figlabel{vislittle}
\vspace{-0.5em}
\end{figure*}
\begin{table*}[t!]
\centering
\renewcommand{\arraystretch}{1.2}
\renewcommand{\tabcolsep}{1.2mm}
\resizebox{\linewidth}{!}{
\begin{tabular}{@{}l|r*{19}{c}|c@{}}
\textbf{VOC 2007 test}        & aero      & bike      & bird      & boat      & bottle     & bus        & car        & cat        & chair      & cow        & table      & dog        & horse      & mbike      & person     & plant      & sheep      & sofa       & train      & tv         & mAP       \\
\hline
R-CNN \pool{5}  &  51.8  &  60.2  &  36.4  &  27.8  &  23.2  &  52.8  &  60.6  &  49.2  &  18.3  &  47.8  &  44.3  &  40.8  &  56.6  &  58.7  &  42.4  &  23.4  &  46.1  &  36.7  &  51.3  &  55.7  &  44.2 \\
R-CNN \fc{6}  &  59.3  &  61.8  &  43.1  &  34.0  &  25.1  &  53.1  &  60.6  &  52.8  &  21.7  &  47.8  &  42.7  &  47.8  &  52.5  &  58.5  &  44.6  &  25.6  &  48.3  &  34.0  &  53.1  &  58.0  &  46.2 \\
R-CNN \fc{7}  &  57.6  &  57.9  &  38.5  &  31.8  &  23.7  &  51.2  &  58.9  &  51.4  &  20.0  &  50.5  &  40.9  &  46.0  &  51.6  &  55.9  &  43.3  &  23.3  &  48.1  &  35.3  &  51.0  &  57.4  &  44.7 \\
\hline
R-CNN FT \pool{5}  &  58.2  &  63.3  &  37.9  &  27.6  &  26.1  &  54.1  &  66.9  &  51.4  &  26.7  &  55.5  &  43.4  &  43.1  &  57.7  &  59.0  &  45.8  &  28.1  &  50.8  &  40.6  &  53.1  &  56.4  &  47.3 \\
R-CNN FT \fc{6}  &  63.5  &  66.0  &  47.9  &  37.7  &  29.9  &  62.5  &  70.2  &  60.2  &  32.0  &  57.9  &  47.0  &  53.5  &  60.1  &  64.2  &  52.2  &  31.3  &  55.0  &  50.0  &  57.7  &  63.0  &  53.1 \\
R-CNN FT \fc{7}  &  64.2  &  69.7  &  50.0  &  41.9  &  32.0  &  62.6  &  71.0  &  60.7  &  32.7  &  58.5  &  46.5  &  56.1  &  60.6  &  66.8  &  54.2  &  31.5  &  52.8  &  48.9  &  57.9  &  64.7  &  54.2 \\
\hline
R-CNN FT \fc{7} BB  &  \bf{68.1}  &  \bf{72.8}  &  \bf{56.8}  &  \bf{43.0}  &  \bf{36.8}  &  \bf{66.3}  &  \bf{74.2}  &  \bf{67.6}  &  \bf{34.4}  &  \bf{63.5}  &  \bf{54.5}  &  \bf{61.2}  &  \bf{69.1}  &  \bf{68.6}  &  \bf{58.7}  &  \bf{33.4}  &  \bf{62.9}  &  \bf{51.1}  &  \bf{62.5}  &  \bf{64.8}  &  \bf{58.5} \\
\hline
\hline
DPM v5 \cite{release5}  &  33.2  &  60.3  &  10.2  &  16.1  &  27.3  &  54.3  &  58.2  &  23.0  &  20.0  &  24.1  &  26.7  &  12.7  &  58.1  &  48.2  &  43.2  &  12.0  &  21.1  &  36.1  &  46.0  &  43.5  &  33.7 \\
DPM ST \cite{lim2013sketch}  &  23.8  &  58.2  &  10.5  &  \phz8.5  &  27.1  &  50.4  &  52.0  &  \phz7.3  &  19.2  &  22.8  &  18.1  &  \phz8.0  &  55.9  &  44.8  &  32.4  &  13.3  &  15.9  &  22.8  &  46.2  &  44.9  &  29.1 \\
DPM HSC \cite{HSC}  &  32.2  &  58.3  &  11.5  &  16.3  &  30.6  &  49.9  &  54.8  &  23.5  &  21.5  &  27.7  &  34.0  &  13.7  &  58.1  &  51.6  &  39.9  &  12.4  &  23.5  &  34.4  &  47.4  &  45.2  &  34.3 \\
\end{tabular}
}
\caption{\textbf{Detection average precision (\%) on VOC 2007 test.}
Rows 1-3 show R-CNN performance without fine-tuning.
Rows 4-6 show results for the CNN pre-trained on ILSVRC 2012
and then fine-tuned (FT) on VOC 2007 trainval.
Row 7 includes a simple bounding-box regression (BB) stage that reduces localization errors (\secref{bboxreg}).
Rows 8-10 present DPM methods as a strong baseline.  
The first uses only HOG, while the next two
use different feature learning approaches to augment or replace HOG. 
}
\tablelabel{voc2007}
%\vspace{-1em}
\end{table*}

\begin{table*}[t!]
\centering
\renewcommand{\arraystretch}{1.2}
\renewcommand{\tabcolsep}{1.2mm}
\resizebox{\linewidth}{!}{
\begin{tabular}{@{}l|r*{19}{c}|c@{}}
\textbf{VOC 2007 test}        & aero      & bike      & bird      & boat      & bottle     & bus        & car        & cat        & chair      & cow        & table      & dog        & horse      & mbike      & person     & plant      & sheep      & sofa       & train      & tv         & mAP       \\
\hline
R-CNN T-Net&  64.2  &  69.7  &  50.0  &  41.9  &  32.0  &  62.6  &  71.0  &  60.7  &  32.7  &  58.5  &  46.5  &  56.1  &  60.6  &  66.8  &  54.2  &  31.5  &  52.8  &  48.9  &  57.9  &  64.7  &  54.2 \\
R-CNN T-Net BB &  68.1  &  72.8  &  56.8  &  43.0  &  36.8  &  66.3  &  74.2  &  67.6  &  34.4  &  63.5  &  54.5  &  61.2  &  69.1  &  68.6  &  58.7  &  33.4  &  62.9  &  51.1  &  62.5  &  64.8  &  58.5 \\
\hline
R-CNN O-Net & 71.6 & 73.5 & 58.1 & 42.2 & 39.4 & 70.7 & 76.0 & 74.5 & 38.7 & 71.0 & 56.9 & 74.5 & 67.9 & 69.6 & 59.3 & \bf{35.7} & 62.1 & 64.0 & 66.5 & \bf{71.2} & 62.2 \\
R-CNN O-Net BB & \bf{73.4} & \bf{77.0} & \bf{63.4} & \bf{45.4} & \bf{44.6} & \bf{75.1} & \bf{78.1} & \bf{79.8} & \bf{40.5} & \bf{73.7} & \bf{62.2} & \bf{79.4} & \bf{78.1} & \bf{73.1} & \bf{64.2} & 35.6 & \bf{66.8} & \bf{67.2} & \bf{70.4} & 71.1 & \bf{66.0} \\
\end{tabular}
}
\caption{\textbf{Detection average precision (\%) on VOC 2007 test for two different CNN architectures.}
The first two rows are results from \tableref{voc2007} using Krizhevsky \etal's architecture (T-Net).
Rows three and four use the recently proposed 16-layer architecture from Simonyan and Zisserman (O-Net) \cite{vggverydeep}.
}
\tablelabel{voc2007oxfordnet}
\vspace{-1em}
\end{table*}


First-layer filters can be visualized directly
and are easy to understand \cite{krizhevsky2012imagenet}. They capture
oriented edges and opponent colors. Understanding the subsequent layers
is more challenging. Zeiler and Fergus present a visually attractive deconvolutional
approach in \cite{zeiler2011adaptive}. We propose a simple (and complementary) non-parametric method that directly shows what the network learned. 

The idea is to single out a particular unit (feature) in the network and use it as if it were an object
detector in its own right. That is, we compute the unit's activations on a
large set of held-out region proposals (about 10 million), sort the proposals from
highest to lowest activation, perform non-maximum suppression, and then display the top-scoring regions. Our method lets the selected unit
``speak for itself'' by showing exactly which inputs it fires on. We avoid
averaging in order to see different visual modes and gain
insight into the invariances computed by the unit.

We visualize units from layer \pool{5}, which is the max-pooled output of
the network's fifth and final convolutional layer.  
The \pool{5} feature map is $6 \times 6 \times 256 = 9216$-dimensional. 
Ignoring boundary effects, each \pool{5} unit has a receptive field of $195
\times 195$ pixels in the original $227 \times 227$ pixel input.
A central \pool{5} unit has a nearly global view, while one near the edge has a smaller, clipped support.
%Layer \pool{5} is the last layer for which units have a compact receptive field, making it easier to show which part of the image is responsible for an activation.
%We can also gain some intuition into the representation
%learned by the next layer, \fc{6}, because it takes multiple weighted combinations of \pool{5} activations.

Each row in \figref{vislittle} displays the top 16 activations for a \pool{5} unit from a CNN that we fine-tuned on VOC 2007 trainval.
Six of the 256 functionally unique units are visualized (\asecref{extravis} includes more).
%The first unit was selected because
%it corresponds to a large positive weight in a person-detector SVM (trained on \pool{5}).
%The remaining units were selected to show a representative sample of what the network learns to represent.
These units were selected to show a representative sample of what the network learns.
In the second row, we see a unit that fires on dog faces and dot arrays.
The unit corresponding to the third row is a red blob detector.
There are also detectors for human faces and more abstract patterns such as text and triangular structures with windows.
The network appears to learn a representation that combines a small number of class-tuned features together with a distributed representation of shape, texture, color, and material properties.
%These visualizations demonstrate the richness and diversity of \pool{5} features. 
The subsequent fully connected layer \fc{6} has the ability to model a large set of compositions of these rich features. 
